\uuid{xwVF}
\titre{Logique}
\auteur{bodin ; croizat ; sacre}
\organisation{exo7}
\licence{CC BY-NC-SA 4.0}
\datecreate{2022-03-08}
\niveau{L1}
\module{Algèbre}
\chapitre{Logique, ensemble, raisonnement}
\sousChapitre{Logique}
\difficulte{3}
\theme{Logique}
\exo7qcmid{hilisit.logique.017}

\contenu{
\begin{question}
\qclassification{100.01}
Quelles sont les assertions vraies ?
\begin{answers}
	\good{$\forall x \in \Rr \; \exists \, y \in \Rr \quad y = e^x$}
	\bad{$\forall y \in \Rr \; \exists \, x \in \Rr \quad y = e^x$}
	\bad{$\forall y \in \Rr \; \exists \, x > 0 \quad y = e^x$}
	\good{$\forall y > 0 \; \exists \, x \in \Rr \quad y = e^x$}
\end{answers}
\begin{explanations}
"$\forall x \in \Rr \; \exists \, y \in \Rr \quad y = e^x$" signifie : pour tout nombre réel $x$, il existe un nombre réel $y$ qui est égal à $e^x$ : c'est vrai ($e^x$ est bien un nombre réel).\\
"$\forall y \in \Rr \; \exists \, x \in \Rr \quad y = e^x$" signifie : pour tout nombre réel $y$, il existe un nombre $x$ tel que $y = e^x$. Ceci est faux puisque si $y<0$, on ne pourra jamais l'écrire comme une exponentielle.\\
"$\forall y \in \Rr \; \exists \, x > 0 \quad y = e^x$" signifie : pour tout nombre réel $y$, il existe un réel strictement positif $x$ tel que $y = e^x$. Ceci est faux : si $y<0$, on ne pourra jamais l'écrire comme une exponentielle. Et si $0 \leq y < 1$, on ne pourra jamais écrire $y$ comme une exponentielle d'un nombre strictement positif.\\
"$\forall y > 0 \; \exists \, x \in \Rr \quad y = e^x$" signifie : pour tout nombre strictement positif $y$, il existe un nombre réel $x$ tel que $y = e^x$. Ceci est vrai : si $y>0$, alors on peut écrire $y = e^{\ln(y)}$. Donc on peut choisir $x = \ln(y)$.
\end{explanations}
\end{question}
}
